% NOTE: In the Word draft this section is flagged by the authors as needing to be substantially
% rewritten to reflect the paper's results and conclusions; it is converted here as drafted.
\section{Concluding Remarks}
\label{sec:conclusion}

The analysis in this paper has not established that all economies will necessarily go through a ``linear'' phase before entering a ``circular'' phase. Whether an economy will follow this course of development will depend on its initial conditions. However, if the early stage of development is characterized by a scarcity of man-made capital relative to natural resources, we have seen that the economy will most likely start out from the linear stage, but at some point it will be induced to move on to a circular stage with some amount of recycling as man-made capital becomes more abundant relative to natural resources. From a social viewpoint the incentive to enter the circular stage will be strengthened in the plausible case where the linear stage involves a decline in the quality of the environment.

The mechanisms that take the economy from the linear to the circular phase could provide part of the explanation for the Environmental Kuznets Curve, according to which pollution increases during the early stage of economic development and declines as the economy matures. In our model the reversal of the tendency for environmental quality to decline is driven by rising natural resource scarcity and an increasing stock of capital, which tend to reduce the use of polluting raw materials and to increase recycling of materials and waste as the relative price of raw material goes up and as the rising supply of capital reduces the opportunity cost of investing in recycling.\footnote{There is a long-standing debate whether an Environmental Kuznets Curve (EKC) actually exists, cf. \textcite{barbier1997} and \textcite{stern2004}. \textcite{brock2010} argue that a Solow growth model with exogenous technical progress in pollution abatement can account for the different EKC patterns observed across countries when one allows for cross-country differences in initial conditions and economic structures.} Our model also indicates that policy makers have an incentive to tighten environmental policy when the growth process causes increasing damage to the environment.

According to our analysis, one important driver of the transition from the linear to the circular economy is a continuing rise in the relative price of exhaustible natural resources. In reality we have not observed an unbroken rising trend in real raw materials prices. \textcite{slade1982} identified a U-shaped time trend in many important resource prices, which could be explained as the net effect of a race between technical progress in extraction technologies and growing scarcity of easily accessible reserve stocks. The more recent research by \textcite{lee2006} suggests that real resource prices have tended to be stationary around deterministic trends with structural breaks. In contrast to the basic Hotelling mechanism emphasized by our model, an important driver of recycling could be technical progress in recycling technologies. Incorporating endogenous technical change into the model would be a natural extension of the present analysis.

Our simplified model has abstracted from many other aspects of recycling that may be important in practice. For example, the treatment of household waste is typically regulated by local governments, and differences in regulation across jurisdictions may hamper the creation of an efficient market for collection, sorting, and recycling of waste, as emphasized by the \textcite{advisoryboard2017} in Denmark. Moreover, the optimal degree of recycling of waste will depend on landfill space constraints \parencite{highfill1997} and on the social value of the waste as a fuel in the production of heat and electricity \parencite{gradus2017}. The choice between home separation of different types of waste by households and post-collection separation of waste such as plastic may also be important for the total costs of recycling, as documented by \textcite{dijkgraaf2017} in the case of the Netherlands. Designing optimal recycling policies thus involves a host of practical issues that are not included in the present analysis.

Furthermore, in practice governments may not have the information and administrative capacity to implement the fine-tuning Pigouvian taxes and subsidies needed to implement the first-best recycling policy, and at any rate the existence of numerous non-environmental market distortions, as well as concerns about the distributional effects of Pigouvian taxes and subsidies, will inevitably force policy makers into a second-best context. Nevertheless, the analysis in this paper suggests that a well-designed package of environmental taxes and subsidies is an important prerequisite for an efficient degree of recycling. At a more basic level, the paper has tried to illustrate that the policy issues studied by proponents of the new circular economy paradigm can be fruitfully analysed by using the tools of conventional environmental economics.
